\oldpage[II]{54}
    à $0$
    (\hyperref[II.2.1.11-fr]{2.1.11}). Dans le cas contraire, on aurait
    $(S_h)_+=0$ pour un $h\neq0$ dans $\mathfrak{J}$, donc
    (\hyperref[II.3.1.1-fr]{3.1.1}) $\varphi^{-1}(D(h))=\varnothing$, et
    $\varphi(X)$ ne serait pas dense dans $Y$ contrairement à l'hypothèse
    (puisque $D(h)\neq\varnothing$, $h$ n'étant pas nilpotent).
\end{enumerate}
\end{proof}

\subsection{Faisceau sur $\operatorname{Proj}(\mathcal{S})$ associé à un
$\mathcal{S}$-Module gradué.}
\label{subsection:II.3.2-fr}

\begin{env}[3.2.1]
\label{II.3.2.1-fr}
Soient $Y$ un préschéma, $\mathcal{S}$ une $\mathcal{O}_Y$-Algèbre graduée
quasi-cohérente à degrés positifs, $\mathcal{M}$ un $\mathcal{S}$-module
gradué quasi-cohérent (sur $(Y,\mathcal{O}_Y)$ ou sur l'espace annelé
$(Y,\mathcal{S})$, ce qui revient au même
(\hyperref[I.9.6.1-fr]{\textbf{I}, 9.6.1})). Avec les notations de
(\hyperref[II.3.1.1-fr]{3.1.1}), désignons par $\widetilde{\mathcal{M}}_U$ le
$\mathcal{O}_{X_U}$-Module quasi-cohérent
$(\Gamma(U,\mathcal{M}))\supertilde$~; pour $U'\subset U$,
$\Gamma(U',\mathcal{M})$ s'identifie canoniquement à
$\Gamma(U,\mathcal{M})\otimes_A A'$ (\hyperref[I.1.6.4-fr]{\textbf{I}, 1.6.4})~; donc
on a $\widetilde{\mathcal{M}}_{U'}
=\rho_{U',U}^*(\widetilde{\mathcal{M}}_U)$
(\hyperref[II.2.8.11-fr]{2.8.11}).
\end{env}

\begin{proposition}[3.2.2]
\label{II.3.2.2-fr}
Il existe sur $\operatorname{Proj}(\mathcal{S})=X$ un
$\mathcal{O}_X$-Module quasi-cohérent et un seul $\widetilde{\mathcal{M}}$ tel
que, pour tout ouvert affine $U$ de $Y$, on ait
$\eta_U^*((\Gamma(U,\mathcal{M}))\supertilde)
=\widetilde{\mathcal{M}}|f^{-1}(U)$
(en désignant par $\eta_U$ l'isomorphisme
$f^{-1}(U)\xrightarrow{\sim}\operatorname{Proj}(\Gamma(U,\mathcal{S}))$, où
$f$ est le morphisme structural $X\to Y$).
\end{proposition}

\begin{proof}
Comme $\rho_{U',U}$ s'identifie au morphisme d'injection
$f^{-1}(U')\to f^{-1}(U)$ (\hyperref[II.3.1.2.1-fr]{3.1.2.1}), la proposition
résulte aussitôt de la relation
$\widetilde{\mathcal{M}}_{U'}=\rho_{U',U}^*(\widetilde{\mathcal{M}}_U)$ et du
principe de recollement des faisceaux (\hyperref[0.3.3.1-fr]{\textbf{0}, 3.3.1}).
\end{proof}

On dit que $\widetilde{\mathcal{M}}$ est le $\mathcal{O}_X$-Module associé au
$\mathcal{S}$-Module gradué quasi-cohérent $\mathcal{M}$.

\begin{proposition}[3.2.3]
\label{II.3.2.3-fr}
Soit $\mathcal{M}$ un $\mathcal{S}$-Module gradué quasi-cohérent, et soit
$f\in\Gamma(Y,\mathcal{S}_d)$ ($d>0$). Si $\xi_f$ est l'isomorphisme canonique
de $X_f$ sur le $Y$-préschéma
$Z_f=\operatorname{Spec}(\mathcal{S}^{(d)}/(f-1)\mathcal{S}^{(d)})$
\textup{(\hyperref[II.3.1.4-fr]{3.1.4})}, $(\xi_f)_*(\widetilde{\mathcal{M}}|X_f)$ est
le $\mathcal{O}_{Z_f}$-Module
$(\mathcal{M}^{(l)}/(f-1)\mathcal{M}^{(d)})\supertilde$
\textup{(\hyperref[II.1.4.3-fr]{1.4.3})}.
\end{proposition}

\begin{proof}
La question étant locale sur $Y$, on est aussitôt ramené à
(\hyperref[II.2.2.5-fr]{2.2.5}), compte tenu de la commutativité du diagramme
(\hyperref[II.2.8.12.1-fr]{2.8.12.1}).
\end{proof}

\begin{proposition}[3.2.4]
\label{II.3.2.4-fr}
Le $\mathcal{O}_X$-Module $\widetilde{\mathcal{M}}$ est un foncteur covariant
additif exact en $\mathcal{M}$, de la catégorie des $\mathcal{S}$-Modules
gradués quasi-cohérents dans celle des $\mathcal{O}_X$-Modules
quasi-cohérents, qui commute aux limites inductives et aux sommes directes.
\end{proposition}

\begin{proof}
La question étant locale sur $Y$, se ramène à
(\hyperref[I.1.3.11-fr]{\textbf{I}, 1.3.11} et
\hyperref[I.1.3.9-fr]{1.3.9}) et à
(\hyperref[II.2.5.4-fr]{2.5.4}).
\end{proof}

En particulier, si $\mathcal{N}$ est un sous-$\mathcal{S}$-Module gradué
quasi-cohérent de $\mathcal{M}$, $\widetilde{\mathcal{N}}$ s'identifie
canoniquement à un sous-$\mathcal{O}_X$-Module quasi-cohérent de
$\widetilde{\mathcal{M}}$~; plus particulièrement, pour tout faisceau gradué
quasi-cohérent $\mathcal{J}$ d'idéaux de $\mathcal{S}$,
$\widetilde{\mathcal{J}}$ est un faisceau quasi-cohérent d'idéaux de
$\mathcal{O}_X$.

Si $\mathcal{M}$ est un $\mathcal{S}$-Module gradué quasi-cohérent et
$\mathcal{I}$ un faisceau quasi-cohérent d'idéaux de $\mathcal{O}_Y$,
$\mathcal{I}\mathcal{M}$ est un sous-$\mathcal{S}$-Module gradué
quasi-cohérent de $\mathcal{M}$ et on a
\[
\label{II.3.2.4.1-fr}
  (\mathcal{I}\mathcal{M})\supertilde
  =\mathcal{I}.\widetilde{\mathcal{M}}
\tag{3.2.4.1}
\]
(le second membre ayant le sens défini en
(\hyperref[0.4.3.5-fr]{\textbf{0}, 4.3.5})). Il suffit en effet de vérifier cette
formule lorsque $Y=\operatorname{Spec}(A)$ est affine,
$\mathcal{S}=\widetilde{S}$, où $S$ est une $A$-algèbre graduée,
$\mathcal{M}=\widetilde{M}$,
